% ============ 11 主齐性空间与 H^1 (原书 pp.76-80) ============
\chapter{\texorpdfstring{主齐性空间与 $H^{1}$}{主齐性空间与 H1}}\label{sec:11}
\index{主齐性空间}
\index{扭子}
\index{1-上链}
\index{Hilbert 定理 90}
\index{Picard 群}
\index{局部自由层}
\index{下降}
\index{gerbe}

在上一节中，我们证明了 $H^{1}(X_{\mathrm{et}}, F)$ 与 $\check{H}^{1}(X_{\mathrm{et}}, F)$ 一致.本节解释如何把 $\check{H}^{1}(X_{\mathrm{et}}, F)$ 解释为 $F$ 的主齐性空间群.由于不难做到(只是定义要稍微小心)，我们处理非交换群层.

\subsection*{第一 \v{C}ech 群的定义}

设 $\mathcal{U} = (U_i \to X)_{i \in I}$ 是 $X$ 的平展覆盖，$G$ 是 $X_{\mathrm{et}}$ 上的群层(不一定是交换的).与上一节一样，我们把 $U_{ij\cdots}$ 写为 $U_i \times_{X} U_j \times_{X} \cdots$.

$\mathcal{U}$ 的取值于 $G$ 的 1-上链是族 $(g_{ij})_{(i,j) \in I \times I}$，其中 $g_{ij} \in G(U_{ij})$，使得
\[
(g_{ij}|_{U_{ijk}}) \cdot (g_{jk}|_{U_{ijk}}) = g_{ik}|_{U_{ijk}}, \qquad \text{所有 } i, j, k.
\]
两个上链 $g$ 与 $g'$ 称为上同调的，记为 $g \sim g'$，如果存在族 $(h_i)_{i \in I}$，$h_i \in G(U_i)$，使得
\[
g'_{ij} = (h_i|_{U_{ij}}) \cdot g_{ij} \cdot (h_j|_{U_{ij}})^{-1}, \qquad \text{所有 } i, j.
\]
1-上链模 $\sim$ 的集合记为 $\check{H}^{1}(X_{\mathrm{et}}, G)$.它一般不是群，但确实有特殊的元素：由对所有 $i, j$ 取 $g_{ij} = 1$ 的 1-上链 $(g_{ij})$ 表示的元素.

群层序列
\[
1 \to G' \to G \to G'' \to 0
\]
称为正合的，如果对所有平展映射 $U \to X$，
\[
1 \to G'(U) \to G(U) \to G''(U)
\]
正合，且 $G \to G''$ 局部满.这样的序列产生集合的序列
\[
1 \to G'(X) \to G(X) \to G''(X) \to \check{H}^{1}(X, G') \to \check{H}^{1}(X, G) \to \check{H}^{1}(X, G'')
\]
它在如下意义下正合：每个箭头的像恰是被下一个箭头映到特殊元素的集合.

\subsection*{主齐性空间}

设 $G$ 是群，$S$ 是 $G$ 右作用的集合.则 $S$ 称为 $G$ 的主齐性空间(或扭子)，如果对某个(从而每个) $s \in S$，映射 $g \mapsto sg \colon G \to S$ 是双射.

设 $G$ 是 $X_{\mathrm{et}}$ 上的群层，$S$ 是 $G$ 右作用的集合层.则 $S$ 称为 $G$ 的主齐性空间，如果
\begin{enumerate}
  \item 存在 $X$ 的平展覆盖 $(U_i \to X)_{i \in I}$，使得对所有 $i$，$S(U_i) \neq \varnothing$；
  \item 对每个平展映射 $U \to X$ 与 $s \in \Gamma(U, S)$，映射 $g \mapsto sg \colon G|_{U} \to S|_{U}$ 是层同构.
\end{enumerate}

主齐性空间 $S$ 称为平凡的，如果它(作为带 $G$ 右作用的层)同构于通过右乘作用于自身的 $G$，或者等价地，如果 $S(X) \neq \varnothing$.公理要求 $S$ 局部同构于平凡主齐性空间.

我们说覆盖 $(U_i \to X)_{i \in I}$ 分裂 $S$，如果 $S(U_i) \neq \varnothing$.

设 $S$ 是 $G$ 的主齐性空间.设 $(U_i \to X)_{i \in I}$ 是分裂 $S$ 的 $X$ 的平展覆盖，对每个 $i$ 选取 $s_i \in S(U_i)$.由条件 (b)，存在唯一的 $g_{ij} \in G(U_{ij})$，使得
\[
(s_i|_{U_{ij}}) \cdot g_{ij} = s_j|_{U_{ij}}.
\]
则 $(g_{ij})_{I \times I}$ 是上链，因为(省略限制记号)
\[
s_i \cdot g_{ij} \cdot g_{jk} = s_k = s_i \cdot g_{ik}.
\]
此外，把 $s_i$ 换成 $s'_i = s_i \cdot h_i$($h_i \in G(U_i)$)给出上同调的上链.因此，$S$ 在 $\check{H}^{1}(\mathcal{U}, G)$ 中定义类 $c(S)$，其中 $\mathcal{U} = (U_i \to X)_{i \in I}$.

\begin{proposition}\label{prop:11.1}
映射 $S \mapsto c(S)$ 定义了被 $\mathcal{U}$ 分裂的 $G$ 的主齐性空间的同构类集合到 $\check{H}^{1}(\mathcal{U}, G)$ 的双射.
\end{proposition}

\begin{proof}
设 $\alpha \colon S \to S'$ 是 $G$-层的同构，选取 $s_i \in S(U_i)$.则 $\alpha(s_i) \in S'$，且(省略限制记号)
\[
s_i \cdot g_{ij} = s_j \Rightarrow \alpha(s_i) \cdot g_{ij} = \alpha(s_j).
\]
因此，由族 $(\alpha(s_i))$ 定义的 1-上链等于由 $(s_i)$ 定义的.这表明 $c(S)$ 只依赖于 $S$ 的同构类，从而 $S \mapsto c(S)$ 确实定义了同构类集合上的映射.

设 $c(S) = c(S')$.则我们可以选取定义相同 1-上链 $(g_{ij})$ 的截面 $s_i \in S(U_i)$ 与 $s'_i \in S'(U_i)$.设存在 $t \in S(X)$.则
\[
t|_{U_i} = s_i \cdot g_i
\]
对唯一 $g_i \in G(U_i)$；由等式 $(t|_{U_i})|_{U_{ij}} = (t|_{U_j})|_{U_{ji}}$，发现
\[
g_{ij} = g_{ij} \cdot (g_j|_{U_{ij}}). \qquad (*)
\]
因为 $S$ 是层，$t \mapsto (g_i)_{i \in I}$ 是 $S(X)$ 到满足 $(*)$ 的族 $(g_i)_{i \in I}$($g_i \in G(U_i)$)的集合上的双射.对 $S'$ 有类似的陈述，从而存在典范双射 $S(X) \to S'(X)$.对任何 $V \to X$，可以对 $V$ 的覆盖 $(U_i \times_{X} V \to V)$ 与元素 $s_i|_{U_i \times_{X} V}$，$s'_i|_{U_i \times_{X} V}$ 应用同样的论证，得到典范双射 $S(V) \to S'(V)$.这些双射的族是 $G$-层的同构 $S \to S'$.

于是，映射是到 $\check{H}^{1}(\mathcal{U}, G)$ 的单射，还剩下证明它满射.

设 $(g_{ij})_{I \times I}$ 是覆盖 $\mathcal{U} = (U_i \to X)$ 的 1-上链.对任何平展映射 $V \to X$，设 $(V_i \to V)$ 是 $V$ 的覆盖，其中 $V_i = U_i \times_{X} V$.定义 $S(V)$ 为满足
\[
(g_i|_{V_{ij}}) = g_{ij} \cdot (g_j|_{V_{ij}})
\]
的族 $(g_i)_{i \in I}$($g_i \in G(V_i)$)的集合.证明这定义 $G$-集合的层且 $c(S)$ 由 $(g_{ij})$ 表示，只涉及常规检验.
\end{proof}

\begin{remark}\label{rem:11.2}
当 $G$ 是由 $X$ 上的群概形 $G$ 定义的层时，人们希望知道每个 $G$ 的主齐性空间 $S$ 都是由概形定义的层.一般地，这可能是下降理论中的困难问题：我们知道对每个 $i \in I$，$S|_{U_i}$ 由 $U_i$ 上的概形表示，即由 $G_{U_i}$ 表示，并且希望知道这蕴含 $S$ 本身由 $X$ 上的概形 $S$ 表示.

已知当例如 $G$ 由 $X$ 上的仿射概形定义时，$S$ 由概形定义.1978 年已知结果的概述见 EC III 4.3---今天知道的也没有多多少.
\end{remark}

\begin{example}\label{ex:11.3}
设 $\underline{G}$ 是由有限群 $G$ 定义的 $X_{\mathrm{et}}$ 上的常值层.$X$ 的带群 $G$ 的 Galois 覆盖是 $\underline{G}$ 的主齐性空间，且每个主齐性空间都来自 Galois 覆盖.当 $X$ 连通时，$X$ 的带群 $G$ 的 Galois 覆盖由连续同态 $\pi_1(X, \bar{x}) \to G$ 分类，其中 $\bar{x} \to X$ 是 $X$ 的任意几何点.于是，对连通的 $X$，有典范同构
\[
H^{1}(X_{\mathrm{et}}, \underline{G}) \cong \operatorname{Hom}_{\mathrm{conts}}(\pi_1(X, \bar{x}), G).
\]
\end{example}

\subsection*{\texorpdfstring{$\check{H}^{1}(X, \mathrm{GL}_n)$ 的解释}{H^1(X, GL_n) 的解释}}

设 $L^{n}(X)$ 是 $X$ 上对 Zariski 拓扑的秩 $n$ 的 $\mathcal{O}_X$-模局部自由层的同构类集合.

\begin{theorem}\label{thm:11.4}
有自然双射
\[
L^{n}(X_{\mathrm{zar}}) \leftrightarrow \check{H}^{1}(X_{\mathrm{zar}}, \mathrm{GL}_n) \leftrightarrow \check{H}^{1}(X_{\mathrm{et}}, \mathrm{GL}_n) \leftrightarrow \check{H}^{1}(X_{\mathrm{fl}}, \mathrm{GL}_n).
\]
\end{theorem}

\begin{proof}
(梗概).设 $\mathcal{M}$ 是 $X_{\mathrm{zar}}$ 上秩 $n$ 的 $\mathcal{O}_X$-模局部自由层.按定义，这意味着存在开覆盖 $X = \bigcup U_i$，使得对所有 $i$，$\mathcal{M}|_{U_i} \cong \mathcal{O}^{n}_{U_i}$.

注意，我们可以取 $U_i$ 仿射，设 $U_i = \operatorname{Spec} A_i$，则 $\mathcal{M}|_{U_i}$ 是与 $A_i$-模 $A^{n}_i$ 相伴的层---特别地，$\mathcal{M}$ 凝聚.

现在对每个 $i$ 选取同构 $\alpha_i \colon \mathcal{M}|_{U_i} \to \mathcal{O}^{n}_{U_i}$.则族
\[
\alpha_{ij} \overset{\mathrm{def}}{=} (\alpha_i|_{U_{ij}}) \circ (\alpha_j|_{U_{ij}})^{-1} \in \mathrm{GL}_n(U_{ij}), \qquad (i, j) \in I \times I
\]
是 1-上链，因为(省略限制记号)
\[
\alpha_{ij} \circ \alpha_{jk} = \alpha_i \circ \alpha^{-1}_{j} \circ \alpha_j \circ \alpha^{-1}_{k} = \alpha_i \circ \alpha^{-1}_{k} = \alpha_{ik} \quad (\text{在 } U_{ijk} \text{ 上}).
\]
此外，若把 $\alpha_i$ 换成 $h_i \circ \alpha_i$，则 $\alpha_{ij}$ 换成
\[
\alpha'_{ij} = h_i \circ \alpha_{ij} \circ h^{-1}_{j},
\]
从而 $(\alpha_{ij})$ 在 $\check{H}^{1}(\mathcal{U}, \mathrm{GL}_n)$($\mathcal{U} = (U_i)$)中的类只依赖于 $\mathcal{M}$.这样，我们得到良定义映射 $L^{n}(X_{\mathrm{zar}}) \to \check{H}^{1}(X_{\mathrm{zar}}, \mathrm{GL}_n)$，并且不难证明它是双射(比较 11.1 的证明).

我们可以类似地定义 $L^{n}(X_{\mathrm{et}})$ 与 $L^{n}(X_{\mathrm{fl}})$，类似的论证表明存在双射 $L^{n}(X_{\tau}) \to \check{H}^{1}(X_{\tau}, \mathrm{GL}_n)$($\tau = \mathrm{et}$ 或 $\mathrm{fl}$).于是，还剩下证明映射 $\mathcal{M} \mapsto \mathcal{M}_{\mathrm{et}}$ 与 $\mathcal{M} \mapsto \mathcal{M}_{\mathrm{fl}}$\footnote{层 $\mathcal{M}_{\mathrm{et}}$ 定义于 \S 6；$\mathcal{M}_{\mathrm{fl}}$ 的定义类似.}给出双射 $L^{n}(X_{\mathrm{zar}}) \to L^{n}(X_{\mathrm{et}})$ 与 $L^{n}(X_{\mathrm{zar}}) \to L^{n}(X_{\mathrm{fl}})$.我们处理平坦情形.必须证明：
\begin{enumerate}
  \item 每个 $X_{\mathrm{fl}}$ 上的 $\mathcal{O}_{X_{\mathrm{fl}}}$-模局部自由层都具有形式 $\mathcal{M}_{\mathrm{fl}}$($\mathcal{M}$ 是 $X_{\mathrm{zar}}$ 上的 $\mathcal{O}_X$-模凝聚层)；
  \item 设 $\mathcal{M}$ 是 $X_{\mathrm{zar}}$ 上的 $\mathcal{O}_X$-模凝聚层；若 $\mathcal{M}_{\mathrm{fl}}$ 局部自由，则 $\mathcal{M}$ 也局部自由；
  \item 对 $L^{n}(X_{\mathrm{zar}})$ 中的 $\mathcal{M}$ 与 $\mathcal{N}$, $\mathcal{M} \cong \mathcal{N}$，当且仅当 $\mathcal{M}_{\mathrm{fl}} \cong \mathcal{N}_{\mathrm{fl}}$.
\end{enumerate}
显然，只需对仿射 $X$ 验证这些论断.

为证明 (b)，需要交换代数的如下结果.

\begin{quote}
\refstepcounter{theorem}\label{item:11.5}
\textbf{11.5} 设 $A$ 是环(交换 Noether 环，一如既往)，$M$ 是有限生成 $A$-模.如下论断等价：
\begin{enumerate}
  \item 对 $A$ 的所有极大理想(从而所有素理想) $\mathfrak{m}$, $A_{\mathfrak{m}}$-模 $M_{\mathfrak{m}}$ 自由；
  \item 存在 $A$ 中的有限元素集 $f_1, \ldots, f_n$，使得 $A = (f_1, \ldots, f_n)$，且对所有 $i$，$M_{f_i}$ 是自由 $A_{f_i}$-模；
  \item $M$ 投射.
\end{enumerate}
见例如 Eisenbud\footnote{\emph{Commutative Algebra}, Springer.} 1995, A3.2.当然，(b) 说的是 $M$ 在 $\operatorname{Spec} A$ 上定义的凝聚层 $\mathcal{M}$(对 Zariski 拓扑)局部自由.
\end{quote}

现在设 $M$ 是有限生成 $A$-模，$A \to B$ 是忠实平坦同态.回忆 $M$ 投射意味着 $\operatorname{Hom}_A(M, \cdot)$ 是正合函子.但对任何 $A$-模 $N$，
\[
B \otimes_{A} \operatorname{Hom}_A(M, N) = \operatorname{Hom}_B(B \otimes_{A} M, B \otimes_{A} N),
\]
从而若 $B \otimes_{A} M$ 投射，则 $M$ 也投射.

现在设 $M$ 定义的 $\mathcal{O}_X$-模凝聚层 $\mathcal{M}$ 在 $X$ 的某个平坦覆盖 $(U_i \xrightarrow{\varphi_i} X)_{i \in I}$ 上变成自由的.以可能扩大 $I$ 为代价，可以假定每个 $U_i$ 仿射.但 $\varphi_i(U_i)$ 在 $X$ 中开，且 $X$ 拟紧，因此可以假定 $I$ 有限.则 $Y = \coprod U_i$ 仿射.于是，存在忠实平坦映射 $A \to B$，使得 $B \otimes_{A} M$ 自由.特别地，$B \otimes_{A} M$ 投射，从而 $M$ 投射.这证明了 (b).

为证明 (a)，我们再次需要交换代数的结果.设 $A \to B$ 是忠实平坦同态.对 $B$-模 $N$，设 $N^{0}$ 与 $N^{1}$ 是分别通过
\[
b \mapsto 1 \otimes b \colon B \to B \otimes_{A} B, \qquad b \mapsto b \otimes 1 \colon B \to B \otimes_{A} B
\]
张量得到的 $B \otimes_{A} B$-模.给出满足某个自然上链条件的同构 $\varphi \colon N^{0} \to N^{1}$\footnote{见 Waterhouse, \emph{Introduction to Affine Group Schemes}, Springer, 1979, 17.2.}就是给出满足 $B \otimes_{A} M = N$ 的 $A$-模 $M$.由于两个映射与 $A \to B$ 的复合相等，$(B \otimes_{A} M)^{0} = (B \otimes_{A} M)^{1}$，从而对 $N = B \otimes_{A} M$，$\varphi$ 可取恒等映射.反过来，给定满足上链条件的 $\varphi$，定义 $M$ 为 $N$ 的如下子模：映射
\[
N \xrightarrow{\text{典范}} N^{1} \quad \text{与} \quad N \xrightarrow{\text{典范}} N^{0} \xrightarrow{\varphi} N^{1}
\]
在它上面一致.

设 $\mathcal{M}$ 是 $X_{\mathrm{fl}}$ 上的 $\mathcal{O}_X$-模层，假定它在某个平坦覆盖上等于 $\mathcal{N}_{\mathrm{fl}}$(如前，可以取为由单个仿射映射 $Y \to X$ 组成).于是 $Y \to X$ 对应于忠实平坦映射 $A \to B$，$\mathcal{M}|_{Y}$ 是由 $B$-模 $N$ 定义的层.复合 $Y \times_{X} Y \rightrightarrows Y \to X$ 相等，因此 $\mathcal{M}$ 到 $Y \times_{X} Y$ 的两个限制相等.恒等映射 $\mathcal{M}|_{Y \times_{X} Y} \to \mathcal{M}|_{Y \times_{X} Y}$ 定义满足上一段上链条件的同构 $\varphi \colon N^{0} \to N^{1}$.因此，对 $(N, \varphi)$，它来自 $A$-模 $M$.最后，由于 $M$ 的构造，$\mathcal{M}$ 是与 $M$ 相伴的层.

这完成了 (a) 的梗概证明，我们省略 (c) 的证明.
\end{proof}

\begin{corollary}\label{cor:11.6}
有典范同构 $H^{1}(X_{\mathrm{et}}, \mathbb{G}_{m}) \cong \operatorname{Pic}(X)$.
\end{corollary}

\begin{proof}
按定义，$\operatorname{Pic}(X)$ 是 $X$ 上对 Zariski 拓扑的秩 1 的 $\mathcal{O}_X$-模局部自由层的同构类群.因此，定理中 $n = 1$ 的情形表明 $\operatorname{Pic}(X) \cong \check{H}^{1}(X_{\mathrm{et}}, \mathbb{G}_{m})$.但由于 $\mathbb{G}_{m}$ 交换，这等于 $H^{1}(X_{\mathrm{et}}, \mathbb{G}_{m})$.
\end{proof}

在 $X = \operatorname{Spec} K$($K$ 是域)的情形，推论说 $H^{1}(X_{\mathrm{et}}, \mathbb{G}_{m}) = 0$.这个论断常称为 Hilbert 定理 90，尽管它是原定理的相当大的推广(原定理见 FT 5.25).该推论也常被称为 Hilbert 定理 90.

\begin{remark}\label{rem:11.7}
事实上，证明推论比证明定理稍容易.对应于连续态射 $\varphi \colon X_{\mathrm{fl}} \to X_{\mathrm{zar}}$，有 Leray 谱序列 $H^{r}(X_{\mathrm{zar}}, R^{s}\varphi_{*}\mathbb{G}_{m}) \Rightarrow H^{r+s}(X_{\mathrm{fl}}, \mathbb{G}_{m})$，其中 $R^{s}\varphi_{*}\mathbb{G}_{m}$ 是与预层 $U \mapsto H^{r}(U_{\mathrm{zar}}, \mathbb{G}_{m})$ 相伴的层(见下一节).该谱序列给出映射 $H^{1}(X_{\mathrm{zar}}, \mathbb{G}_{m}) \to H^{1}(X_{\mathrm{fl}}, \mathbb{G}_{m})$，若 $R^{1}\varphi_{*}\mathbb{G}_{m} = 0$ 则它是同构.要检验这一点，只需证明 $R^{1}\varphi_{*}\mathbb{G}_{m}$ 的茎为零，而这归结为证明：当 $Y = \operatorname{Spec} \mathcal{O}_{X, x}$($x \in X$ 适当)时，$H^{1}(Y_{\mathrm{fl}}, \mathbb{G}_{m}) = 0$.于是，我们只需在上面的证明中取 $A$ 为局部环且 $n = 1$ 的情形给出论证.这稍容易一些.
\end{remark}

\subsection*{\texorpdfstring{非交换 $H^{2}$}{非交换 H2}}

可以对任何群层 $G$(不一定是交换的)定义 $H^{2}(X_{\mathrm{et}}, G)$(事实上，对更一般的对象：束(band，法语 lien)).该集合分类被 $G$ 束缚的 gerbe(法语 gerbes)的等价类.这个理论的简短但并非完全可靠\footnote{p.223 第一段中的一些论断只对交换的 $F$ 成立.}的概述，见 Deligne et al, \emph{Hodge Cycles, Motives, and Shimura Varieties}, Lecture Notes in Math. 900, Springer, 1982, pp 220--228.
